% SPDX-FileCopyrightText: 2026 Yiyu Zhou <yiyu@yiyuzhou.io>
%
% SPDX-License-Identifier: GPL-3.0-or-later

\chapter{Lights and network cues}
\label{ch:lights}

Lighting consoles, dimmers, and nodes take levels over a network as sACN,
the open standard ANSI E1.31, or as Art-Net.  cinq sends both, so light cues
can set the lights along with the sound.  Network cues send Open Sound
Control (OSC) messages, which can fire a cue on a lighting console or in
another program, and cinq takes OSC commands in return, so that a tablet or a
console can fire its cues.  This chapter needs no lighting hardware, since the
window shows every level that cinq sends.

Press \key{Control+l} and open \texttt{lights.cinq.toml} in your copy of the
demo.

\section{The light patch}

Each light is an instrument with a name and a DMX address.  The address is
the channel that sets the intensity of the instrument.  The instruments form
the light patch.

\begin{enumerate}
\item Click the \ui{Light Patch} tab of the inspector.  It lists the three
  instruments of the show, as in Figure~\ref{fig:patch}.
\item Press \ui{GO} twice.  Cue~1 brings the front light to 30\% at once.
  Cue~2 fades it to 70\% and the back light to 40\% over three seconds.  The
  \ui{Now} column follows each level.
\end{enumerate}

\begin{figure}[H]
  \centering
  \includegraphics[width=\textwidth]{screenshots/patch.png}
  \caption{\ui{Light Patch} tab after cue~2}
  \label{fig:patch}
\end{figure}

Universes count from 1, as in sACN, and each holds 512 addresses.  The
\ui{Address} column shows an address in universe~1 as the address alone.  In
another universe, it shows the universe, a slash, and the address, as
\texttt{2/1} for the lightning.  \ui{Add instrument} adds an instrument at
the address after the last one.  \ui{Remove} removes an instrument from the
patch, and no cue sets it afterward.  A light cue that still names it reads
\texttt{UNPATCHED} in yellow in its state, and one that names no instrument
reads \texttt{NO LEVELS}.  The inspector says why in yellow, and the sidebar
lists the cue among the warnings.  On the \ui{Levels} tab of an unpatched cue,
the missing instrument reads \texttt{99 (unpatched)}, and clearing its field
puts the cue right.  With the keys, \key{Tab} reaches each of these buttons, in
the order of the window, and \key{Enter} presses it.  The table scrolls to show
the field or \ui{Remove} button that \key{Tab} reached.

\section{Light cues}

A light cue fades the instruments that it names to its own levels, in
percent, over its duration.  The other instruments keep their levels until a
later light cue names them, as in QLab.

\begin{enumerate}
\item Select cue~2, Stage up, and click the \ui{Levels} tab.  It lists every
  instrument with its level in the cue and the level that it has on stage
  now, as in Figure~\ref{fig:light}.  An empty field, as for the lightning,
  leaves the instrument unchanged.
\item Press \ui{GO}\@.  Cue~3 flashes the lightning.  It auto-continues to
  cue~4, which fades the lightning out in half a second, and cue~4
  auto-continues at once to cue~5, which plays the thunder.
\end{enumerate}

\begin{figure}[H]
  \centering
  \includegraphics[width=\textwidth]{screenshots/light.png}
  \caption{\ui{Levels} tab of a light cue}
  \label{fig:light}
\end{figure}

Pausing a light cue pauses its fade.  Stopping it, as Panic does, leaves the
lights where the fade has brought them, so Panic never blacks out the stage.
For a blackout, use a light cue that takes every instrument to 0\%, such as
cue~7.

\section{Seeing the output}

cinq sends each universe of the patch as sACN and as Art-Net.  It sends them
on every network of the computer and to the computer itself, whenever a level
changes and again every second.  An sACN monitor, such as sACNView, shows
universes 1 and 2 as the cues run.  cinq also announces its universes every
ten seconds, so the monitor lists the source \texttt{cinq} without further
setup.  On Darwin, allow cinq to find devices on local networks when it asks.

A receiver merges the sACN of cinq with that of any other source at the same
priority that drives the same universe, such as a lighting console, and takes
the higher level of each address.  Test on a universe that nothing else uses.
Drive the universe of a real rig only with the agreement of the lighting
designer.

\section{Network cues}

A network cue sends one OSC message when it starts.  For TCP, the destination
is \texttt{host:port}.  For UDP, it is \texttt{udp://host:port}.  The message
is an address that starts with a slash, followed by arguments.  Each argument
is a number or text.  Double quotes keep text with spaces whole.

\begin{enumerate}
\item Select cue~6 and click the \ui{Basics} tab.  The \ui{Destination} and
  \ui{Message} rows hold the destination, UDP port 9000 on this computer, and
  the message, as in Figure~\ref{fig:network}.
\item In a terminal, start a program that prints the OSC messages that it
  receives on that port, such as \texttt{oscdump 9000} from liblo.
\item Press \ui{GO}\@.  After the thunder, cue~6 is on standby, so GO fires
  it.  The program prints the message, with \texttt{Act 1} and \texttt{3} as
  its arguments.
\end{enumerate}

\begin{figure}[H]
  \centering
  \includegraphics[width=\textwidth]{screenshots/network.png}
  \caption{\ui{Basics} tab of a network cue}
  \label{fig:network}
\end{figure}

Each console or program documents the messages that it accepts and the port
that it listens on.  For example, the Eos consoles of ETC take OSC over TCP
with SLIP framing on port~3037, and \texttt{/eos/cue/1/5/fire} fires cue~5
of cue list~1.

A message that cannot be sent, as when no program accepts its TCP
connection, appears among the warnings in the sidebar, for the first cue
that names the destination, and the \texttt{State} of every cue that names
the same destination reads \texttt{SEND FAILED}.  Pressing the warning selects
the cue, whose inspector shows the destination.  cinq drops a message that it
cannot send within a second of its cue, since a console cue that fires late
can do more harm than one that does not fire.  A UDP message fails in the same
way when its host cannot be found, but one that goes out and reaches no program
reports no failure, since nothing answers it.

\section{Commands from the network}

A tablet, a console, or another program can fire cues over the network with
OSC commands.  Add this line to \texttt{config.toml}, at the top, above any
line that starts with \texttt{[}, since a setting below a table belongs to the
table.  The help page links to the file, as Chapter~\ref{ch:keys} tells.  cinq
rereads the file when it is saved:

\begin{verbatim}
osc_listen = "127.0.0.1:53000"
\end{verbatim}

The address \texttt{127.0.0.1} lets programs on this computer alone send
commands.  Anyone who can reach the address and port can fire cues, so for
another computer give the address of the interface that the show network
uses, such as \texttt{192.168.1.20:53000}, and keep the port from other
networks.  Without the line, cinq listens to nothing.

The commands are those of QLab, so remote controls made for QLab work as they
are.  \texttt{/go} and \texttt{/panic} act as the keys do, and \texttt{/pause}
and \texttt{/resume} act on every running cue.  A cue goes by its number as the
cue list shows it: \texttt{/cue/5/start}, \texttt{/cue/5/stop},
\texttt{/cue/5/pause}, \texttt{/cue/5/resume}, and \texttt{/playhead/5}, which
moves standby to cue~5.  A program such as \texttt{oscsend} from liblo sends
one, to the address of the setting:

\begin{verbatim}
oscsend 127.0.0.1 53000 /go
\end{verbatim}

Without liblo, bash sends the same message.  An OSC string ends in a zero byte
and is padded with zero bytes to a multiple of four bytes, so that \texttt{/go}
takes one zero byte, and the comma that lists no arguments takes three:

\begin{verbatim}
printf "/go\0,\0\0\0" > /dev/udp/127.0.0.1/53000
\end{verbatim}

cinq fires the cue that is on standby, as \ui{GO} does.  A remote control
whose button sends a 1 when pressed and a 0 when released fires once, since a
message whose only argument is 0 does nothing.  The sidebar warns when the
address cannot be used, as when another program holds the port.  A command
that does nothing, such as one for a cue that does not exist, shows in the
sidebar among the warnings, with the address that it came from, as in
\texttt{OSC from 127.0.0.1: no cue is numbered "99"}, and stays until a
command acts.  It goes to standard error as well, which a launcher does not
show.  The manual page \texttt{cinq(5)} describes the setting in full.

\section*{Exercise}

Add an instrument with \ui{Add instrument}.  Then choose \ui{Light} in the
\ui{New cue} drop-down list and press \ui{Add cue}.  Set the \ui{Duration} of
the new cue to 2 seconds on the \ui{Basics} tab, which holds the fade time of
a light cue, and on the \ui{Levels} tab of the cue, set the
level of the new instrument to 50.  Press \ui{Standby selected}, then
\ui{GO}\@.  Watch the level in the \ui{Now} column or in your sACN monitor.
